% !TEX root = ../main.tex
\section{Problema 2: Kakuro}

\subsection{Descripción del Problema y Abstracción CSP}
% Descripción del juego tipo crucigrama numérico de sumas cruzadas.
% Explicar el objetivo: rellenar casillas con dígitos 1 al 9 respetando sumas y sin repetir dígitos en cada tramo.

\subsubsection{Parámetros y Datos de Entrada}
% Especificar cómo se representan las casillas blancas y las pistas de sumas horizontales y verticales.

\subsubsection{Variables de Decisión y Dominios}
% Definición precisa de las variables (ej. valor en cada casilla blanca).
% Dominio: {1, ..., 9}.

\subsubsection{Restricciones del Modelo}
% Formalizar en lenguaje matemático:
% 1. Suma de cada serie (horizontal/vertical) igual a la pista dada.
% 2. Todos los dígitos en cada serie deben ser diferentes (alldifferent).

\subsection{Detalles de Implementación (\texttt{kakuro.mzn})}
% Describir aspectos relevantes de la implementación sin copiar el código completo.

\subsubsection{Variantes en el Modelado de Restricciones}
% Comparar enfoques de modelado y propagación para sumas y alldifferent.

\subsubsection{Restricciones Redundantes y Rompimiento de Simetrías}
% Explicar si aplican restricciones redundantes (ej. límites teóricos min/max de suma según longitud) 
% o justificar si no hay simetrías.

\subsection{Estrategias de Búsqueda y Distribución}
% Describir y comparar al menos dos anotaciones de búsqueda (selección de variable y valor).

\subsection{Pruebas y Análisis Experimental}
% Presentar la batería de pruebas y las tablas comparativas con estadísticas de MiniZinc.

\begin{table}[htbp]
    \centering
    \begin{tabular}{lcccc}
        \toprule
        Instancia & Estrategia & Tiempo (s) & Nodos explorados & Nodos fallidos \\
        \midrule
        Kakuro 1 & Default & -- & -- & -- \\
        Kakuro 1 & Custom & -- & -- & -- \\
        \bottomrule
    \end{tabular}
    \caption{Comparación experimental del modelo de Kakuro.}
    \label{tab:kakuro_resultados}
\end{table}

\subsection{Conclusiones del Ejercicio}
% Análisis sobre la eficiencia de las búsquedas y el comportamiento del árbol de búsqueda.
